%%%%%%%%%%%%%%%%%%%%%%%%%%%%%%%%%%%%%%%%%%%%%%%%%%%%%%%%%%%%%%%%%%
\section{Construction of the manufactured FSI solution}
\label{app:mms}
%%%%%%%%%%%%%%%%%%%%%%%%%%%%%%%%%%%%%%%%%%%%%%%%%%%%%%%%%%%%%%%%%%
% Source: paper-data/mms_fsi/paper/mms_construction.tex (24 Nov 2026),
% adapted; every formula follows paper-data/mms_fsi/derive/mms_fsi.py, with
% the checks in paper-data/mms_fsi/derive/generated*.{derive,check}.log.
% Notation of Section 2: v_F, D_S, sigma_S, P_S, kinematic p_F,
% OpenFOAM gradient convention (grad a)_ij = d a_j / d x_i,
% F = I + (grad_0 D_S)^T.

This appendix gives the closed form of the manufactured solution of Section~\ref{sec:mms_fsi}, the design decisions behind it, its boundary data and error measures, and the checks of the construction and of the generated code.
The derivation script, the code generator, the checks, the case libraries and the case generator are in \texttt{paper-data/mms\_fsi/} (\texttt{derive/} and \texttt{case/}).
The aim is a solution that (i) satisfies the kinematic and dynamic interface conditions~\eqref{eqn:kinematic_interface} exactly, so that the coupling is tested without an interface source; (ii) has a finite-strain solid and an interface that moves normal and tangential to itself, so that the fluid domain and its mesh deform; (iii) has fluid and solid fields of physically consistent size, in particular a fluid wall shear of the order of $\mu_F$ times a velocity gradient; and (iv) is closed form, so that it can be evaluated to round-off at any point.

The fluid occupies $0 < x < L$, $0 < y < h(x,t)$, with an inlet at $x = 0$, an outlet at $x = L$, a no-slip wall at $y = 0$ and the interface $y = h(x,t)$.
The solid occupies $0 < X < L$, $H < Y < H + T$ in its initial configuration, with the interface at $Y = H$, a clamped top at $Y = H + T$, and ends at $X = 0$ and $X = L$.
The parameters are given in Table~\ref{tab:mms_params}; all quantities are in consistent units, and the density ratio is one, so added-mass effects are strong.

\begin{table}[htbp]
	\centering
	\small
	\caption{Manufactured FSI solution: parameters (exact rationals in the derivation). The flat variant uses the same values, with the interface of Equation~\eqref{eq:mms_flat}.}
	\label{tab:mms_params}
	\begin{tabular}{lll}
		\toprule
		Parameter & Symbol & Value \\
		\midrule
		Channel length, height; wall thickness & $L$, $H$; $T$ & $2$, $1$; $1/2$ \\
		Angular frequency (period 1) & $\omega$ & $2\pi$ \\
		Normal interface amplitude; travelling-wave fraction & & $1/10$; $1/2$ \\
		Tangential stretch amplitude & & $1/10$ \\
		Through-thickness shape factor & & $1/4$ \\
		Interface shear amplitude & & $1/50$ \\
		Outlet flux amplitude; vortex amplitude & & $1$; $8$ \\
		Free pressure amplitude & & $1/2$ \\
		Fluid density, dynamic viscosity & $\rho_F$, $\mu_F$ & $1$, $1/100$ \\
		Solid density, Young's modulus, Poisson's ratio & $\rho_{S,0}$, $E_S$, $\nu_S$ & $1$, $10$, $3/10$ \\
		\bottomrule
	\end{tabular}
\end{table}

\paragraph{Time factors.}
The motion is built from the time factors
\begin{equation}
	r(t) = \tfrac{1}{2}\left(1 - \cos\omega t\right), \qquad r_2(t) = \sin^2\omega t,
	\label{eq:mms_r}
\end{equation}
which vanish with their first derivatives at $t = 0$.
The displacement, the velocities and the pressure are therefore zero initially (the sources and the solid acceleration are not), and the calculation starts from rest on an undeformed mesh.

\paragraph{Interface motion, built from the solid side.}
The kinematic condition holds at the current position of each material interface point, $\bb{x}^\Gamma(X,t) = (x^\Gamma, y^\Gamma)$, whereas the fluid is written in Eulerian coordinates.
The construction therefore starts from the solid, and uses an interface parametrisation $X \mapsto x^\Gamma$ with a closed-form inverse.
The tangential motion is a M\"obius map of $[0,L]$ onto itself,
\begin{equation}
	x^\Gamma(X,t) = \frac{L\,q}{q + \kappa(1-q)}, \qquad
	\Xi(x,t) = \frac{L\,\kappa\,\hat{x}}{1 - \hat{x} + \kappa\,\hat{x}}, \qquad
	q = \frac{X}{L}, \quad \hat{x} = \frac{x}{L}, \quad \kappa(t) = 1 + \tfrac{1}{10}\,r_2(t),
	\label{eq:mms_mobius}
\end{equation}
where $\Xi$ is the inverse, $\Xi(x^\Gamma(X,t),t) = X$.
It fixes both end points and stretches the interface tangentially, non-uniformly, by up to $10\%$.
The normal motion is prescribed in Eulerian form, as the channel height
\begin{equation}
	h(x,t) = H + \hat\eta(x,t), \qquad
	\hat\eta = \tfrac{1}{10}\,\sin^2\!\left(\frac{\pi x}{L}\right) r(t)\left[1 + \tfrac{1}{2}\sin\!\left(\frac{2\pi x}{L} - \omega t\right)\right],
	\label{eq:mms_height}
\end{equation}
a pulsing wave with a travelling part.
The interface displacement is $\xi = x^\Gamma - X$ tangentially and $\eta(X,t) = \hat\eta(x^\Gamma(X,t),t)$ normally.
Prescribing the height in Lagrangian form would make the swept area $\int h\,\mathrm{d}x$, which the fluid flux below needs, non-elementary; in Eulerian form it is elementary.
The factor $\sin^2(\pi x/L)$ makes $\hat\eta$ and its slope vanish at the ends, where the solid is clamped, which the shear construction below needs.
In the \emph{flat} variant the interface stays flat and moves only vertically,
\begin{equation}
	\hat\eta(t) = \tfrac{1}{10}\,r(t)\left[1 + \tfrac12\sin\omega t\right], \qquad \kappa = 1,
	\label{eq:mms_flat}
\end{equation}
so that the fluid mesh stays orthogonal at the interface; the solid ends then carry the exact displacement instead of a clamp, and the other definitions are unchanged.

\paragraph{Solid displacement and the neo-Hookean shear identity.}
With $\hat{y} = (Y - H)/T$ and $g = (1 - \hat{y})(1 + \tfrac14\hat{y})$, so that $g = 1$ on the interface and $g = 0$ on the top,
\begin{equation}
	D_{S,x} = \xi\,g + (Y - H)(1 - \hat{y})^2\,(m_0 + m_1), \qquad
	D_{S,y} = \eta\,g ,
	\label{eq:mms_DS}
\end{equation}
where $m_0$ and $m_1$ change only the through-thickness shear on the interface.
$\bb{D}_S$ vanishes on the top and, in the full variant, on both ends, where $\xi$, $\eta$, $\eta_{,X}$, $m_0$ and $m_1$ vanish; the interface is then the only solid boundary with non-zero data.
The stress is the neo-Hookean law~\eqref{eqn:neo_hookean} in plane strain ($F_{zz} = 1$), with $\bb{F} = \bb{I} + (\bb{\nabla}_0\bb{D}_S)^T$ and $J = \det\bb{F}$ as in Section~\ref{sec:solid_model}.

For an arbitrary displacement the interface shear traction is of order $\mu_S$ times the strain, i.e.\ of order one here; balancing it with $\mu_F = 10^{-2}$ needs fluid shear rates of order $10^2$, and gave fluid velocities of about 30 in a first attempt.
For the law~\eqref{eqn:neo_hookean}, let $\bb{a} = \partial\bb{x}/\partial X$ and $\bb{b} = \partial\bb{x}/\partial Y$ be the deformed material tangent and through-thickness fibre on the interface.
The volumetric and the isotropic terms give no shear, and the deformed normal is parallel to $\bb{F}^{-T}\bb{N}$, with $|\bb{F}^{-T}\bb{N}| = |\bb{a}|/J$ in two dimensions.
The shear traction on the interface is therefore
\begin{equation}
	\bb{\tau}\cdot\bb{\sigma}_S\cdot\hat{\bb{n}}_F = \pm\,\mu_S\,J^{-2/3}\,\frac{\bb{a}\cdot\bb{b}}{|\bb{a}|^2},
	\label{eq:mms_shear}
\end{equation}
where $\bb{\tau} = \bb{a}/|\bb{a}|$ is the unit tangent of the deformed interface (with the solid normal $\hat{\bb{n}}_S = -\hat{\bb{n}}_F$ the sign is reversed).
The shear vanishes if and only if the deformed fibre $\bb{b}$ stays orthogonal to the deformed interface.
On $Y = H$, $\bb{a} = (1 + \xi_{,X},\ \eta_{,X})$ and $\bb{b} = (\xi g_0' + m_0 + m_1,\ 1 + \eta g_0')$, with $g_0' = g'(H) = -3/(4T)$.
The choice
\begin{equation}
	m_0 = -\frac{\eta_{,X}\left(1 + \eta\,g_0'\right)}{1 + \xi_{,X}} - \xi\,g_0', \qquad
	m_1 = \tfrac{1}{50}\,\sin\!\left(\frac{2\pi X}{L}\right) r(t)
	\label{eq:mms_m}
\end{equation}
makes $\bb{a}\cdot\bb{b} = 0$ for $m_1 = 0$, and $m_1$ then adds a small, smooth shear of known size: at most $0.08$ in magnitude, against a normal traction of up to $3.0$ (maxima over a $401\times401$ grid in $X$ and $t$), so that the fluid wall shear stays of the order of $\mu_F$ times a velocity gradient.
The identity is specific to this law: for Saint Venant--Kirchhoff the interface shear is quadratic in $\bb{b}$, and another device would be needed.
Over a sampling grid the solid reaches a Green--Lagrange strain (Frobenius norm) of $0.26$, a rotation of $0.23$\,rad and $0.69 \le J \le 1.09$.

\paragraph{Divergence-free fluid velocity.}
The fluid velocity is the curl of a stream function in the stretched coordinate $s = y/h(x,t) \in [0,1]$,
\begin{equation}
	\bb{v}_F = \left(\frac{\partial\psi}{\partial y},\, -\frac{\partial\psi}{\partial x}\right), \qquad
	\psi = Q\,B(s) + h\,V_x\,C(s) + s^2(1-s)^2\,\Phi(x,t) + s^2(1-s)^3\,\Omega(x,s,t),
	\label{eq:mms_psi}
\end{equation}
so that $\bb{\nabla}\cdot\bb{v}_F = 0$ identically and the pressure-correction algorithms need no mass source.
Here $\bb{V}(x,t) = \partial\bb{D}_S/\partial t\,(\Xi(x,t), H, t)$ is the interface velocity at the Eulerian abscissa, and $B = 3s^2 - 2s^3$ and $C = s^3 - s^2$ are Hermite functions.
Every term vanishes with its $s$-derivative at $s = 0$ (no slip).
At $s = 1$, $\psi = Q$ and $\partial\psi/\partial y = V_x$, so the tangential velocity matches, and $v_{F,y} = -Q_{,x} + h_{,x}V_x$.
With $Q_{,x} = -h_{,t}$ this is $h_{,t} + h_{,x}V_x = V_y$, the rate of change of the height following the sliding interface point, so $\bb{v}_F = \bb{V}$ on the interface.
The volume flux through the section at $x$,
\begin{equation}
	Q(x,t) = r(t) + \frac{\partial}{\partial t}\int_x^L \hat\eta(x',t)\,\mathrm{d}x',
	\label{eq:mms_Q}
\end{equation}
prescribes a non-negative outlet flux $Q(L,t) = r(t)$; the inlet absorbs the change of volume of the channel, and the outlet, where the velocity gradient is prescribed, never sees backflow, which would make it unstable.
The last term of~\eqref{eq:mms_psi} is a vortex train travelling downstream,
\begin{equation}
	\Omega = 8\,r(t)\,\sin^2\!\left(\frac{\pi x}{L}\right)\sin\!\left(\frac{2\pi x}{L} - \omega t\right)\left(1 + \tfrac{1}{2}\cos\pi s\right).
\end{equation}
It is faded out at both ends so that it does not reverse the outlet flow, and its weight $s^2(1-s)^3$ vanishes with two derivatives at $s = 1$, so it changes neither the interface velocity nor the interface traction.
The peak speed is about $2.1$, and the Reynolds number $\rho_F |\bb{v}_F|_{\max} H/\mu_F$ is about $200$.

\paragraph{Traction continuity without an interface source.}
On $y = h$, the unit normal and tangent are $\hat{\bb{n}}_F = (-h_{,x}, 1)/\sqrt{1 + h_{,x}^2}$ and $\bb{\tau} = (1, h_{,x})/\sqrt{1 + h_{,x}^2}$, and $\hat{\bb{n}}_F = -\hat{\bb{n}}_S$, where $\hat{\bb{n}}_S = \bb{F}^{-T}\bb{N}/|\bb{F}^{-T}\bb{N}|$ is the outward normal of the deformed solid (Nanson's formula) and $\bb{N} = (0,-1)$.
The fluid traction $\rho_F\,\hat{\bb{n}}_F\cdot\bb{\sigma}_F^*$ is linear in two functions that are still free on the interface: $\Phi$ and the interface pressure.
The term $s^2(1-s)^2\Phi$ has zero value and slope at $s = 1$, so $\Phi$ leaves the interface velocity unchanged and enters the interface shear rate only through its second $s$-derivative.
The tangential balance is therefore algebraic in $\Phi$,
\begin{equation}
	c_0 + c_1\,\Phi = \bb{\tau}\cdot\bb{\sigma}_S\cdot\hat{\bb{n}}_F,
	\qquad c_1 = \frac{2\mu_F\left(1 + h_{,x}^2\right)}{h^2},
	\label{eq:mms_Phi}
\end{equation}
where $c_0$ is the fluid shear traction of the other terms of~\eqref{eq:mms_psi}, and $\bb{\sigma}_S$ is evaluated at $(\Xi(x,t), H, t)$; this gives $\Phi$ in closed form.
The normal balance gives the interface pressure, and a term that vanishes on the interface completes the pressure field:
\begin{equation}
	\rho_F\,p_F = \rho_F\,p_\Gamma + (1 - s)\,\tfrac{1}{2}\,r(t)\cos\!\left(\frac{\pi x}{L}\right), \qquad
	\rho_F\,p_\Gamma = 2\mu_F\,\hat{\bb{n}}_F\cdot\bb{d}\cdot\hat{\bb{n}}_F - \hat{\bb{n}}_F\cdot\bb{\sigma}_S\cdot\hat{\bb{n}}_F,
	\label{eq:mms_p}
\end{equation}
with $\bb{d}$ the rate of strain of $\bb{v}_F$.
Both conditions~\eqref{eqn:kinematic_interface} then hold exactly, and no interface source is needed.
An interface source would have required a new term, implemented consistently, in each interface treatment (the quasi-monolithic traction transfer, the Neumann condition of the Dirichlet--Neumann coupling and the Robin condition), and would have bypassed part of the traction transfer under test.

\paragraph{Sources.}
No mass source is needed, since $\bb{\nabla}\cdot\bb{v}_F = 0$ identically, and no interface source, since both interface conditions hold exactly.
The fluid and solid momentum equations~\eqref{eqn:fluid_momentum_ale} and~\eqref{eqn:solid_momentum_TL} carry the body forces
\begin{equation}
	\bb{b}_F = \frac{\partial\bb{v}_F}{\partial t} + \bb{v}_F\cdot\bb{\nabla}\bb{v}_F - \bb{\nabla}\cdot\bb{\sigma}_F^*, \qquad
	\rho_{S,0}\,\bb{b}_S = \rho_{S,0}\,\frac{\partial^2\bb{D}_S}{\partial t^2} - \bb{\nabla}_0\cdot\bb{P}_S^T ,
	\label{eq:mms_sources}
\end{equation}
the residuals of the exact fields.
Here $\bb{v}_F$ is the velocity of Equation~\eqref{eq:mms_psi} and $p_F$ the kinematic pressure of Equation~\eqref{eq:mms_p}; $\bb{\sigma}_F^* = \nu\left[\bb{\nabla}\bb{v}_F + (\bb{\nabla}\bb{v}_F)^T\right] - p_F\bb{I}$, with $\nu = \mu_F/\rho_F$, so that $\bb{\nabla}\cdot\bb{\sigma}_F^* = \nu\nabla^2\bb{v}_F - \bb{\nabla}p_F$; $\bb{D}_S$ is the displacement of Equation~\eqref{eq:mms_DS}, $\bb{P}_S = J\bb{\sigma}_S\bb{F}^{-T}$ with $\bb{\sigma}_S$ from Equation~\eqref{eqn:neo_hookean}, and $\rho_{S,0}$ the initial solid density.
The time derivative of $\bb{v}_F$ is taken at fixed $\bb{x}$ and that of $\bb{D}_S$ at fixed $\bb{X}$.
In Cartesian components, $(\bb{v}_F\cdot\bb{\nabla}\bb{v}_F)_i = v_j\,\partial v_i/\partial x_j$ and $(\bb{\nabla}_0\cdot\bb{P}_S^T)_i = \partial P_{iJ}/\partial X_J$, summed over repeated indices.
After differentiation the closed forms run to thousands of terms and are not printed: they are generated from these definitions by the derivation script \texttt{paper-data/mms\_fsi/derive/mms\_fsi.py}, and the generated C++ evaluation code, \texttt{derive/generated/mmsFsiExact.H} (\texttt{generated\_flat/} for the flat variant), gives $\bb{v}_F$, $p_F$, $\bb{D}_S$, $\bb{b}_F$ and $\rho_{S,0}\bb{b}_S$ at any point and time, so that the case can be reproduced in another code without the symbolic pipeline.
$\bb{b}_F$ is Eulerian, and applies unchanged on the moving mesh, since $\partial_t\bb{v}_F|_{\chi} + (\bb{v}_F - \bb{v}^\Gamma)\cdot\bb{\nabla}\bb{v}_F = \partial_t\bb{v}_F|_{\bb{x}} + \bb{v}_F\cdot\bb{\nabla}\bb{v}_F$ for any velocity $\bb{v}^\Gamma$ of the mesh points, with $\chi$ the mesh coordinates.
It is per unit mass, as the fluid code uses the kinematic pressure, and is evaluated at the current cell centres at the new time level; $\bb{b}_S$ is evaluated at the initial cell centres, and the solid source per unit initial volume, $\rho_{S,0}\bb{b}_S$, is what the code applies.
Both are of the size of the physical terms (maxima over a sampling grid, full variant): $|\bb{b}_F| \approx 10$ against $|\partial\bb{v}_F/\partial t| \approx 9$, $|\bb{v}_F\cdot\bb{\nabla}\bb{v}_F| \approx 6$ and $|\bb{\nabla}p_F| \approx 8$; and $|\rho_{S,0}\bb{b}_S| \approx 10$ against $|\rho_{S,0}\,\partial^2\bb{D}_S/\partial t^2| \approx 5$ and $|\bb{\nabla}_0\cdot\bb{P}_S^T| \approx 11$.

\paragraph{Boundary data.}
Every prescribed boundary value and gradient is the exact one (Table~\ref{tab:mms_bc}); the interface data, and the momentum-balance and flux-consistent pressure gradients, are computed by the solver from the numerical fields.

\begin{table}[htbp]
	\centering
	\small
	\caption{Boundary conditions of the manufactured FSI solution. ``Exact'' means the manufactured value or normal gradient evaluated at the face centre at the new time level; QM: quasi-monolithic fluid.}
	\label{tab:mms_bc}
	\begin{tabular}{lp{3.2cm}p{6.2cm}}
		\toprule
		Boundary & Velocity / displacement & Pressure \\
		\midrule
		Inlet, $x = 0$ & exact $\bb{v}_F$ & QM: exact $\partial p_F/\partial n$; partitioned: flux-consistent gradient \\
		Outlet, $x = L$ & exact $\partial\bb{v}_F/\partial n$ & exact $p_F$ (sets the level); partitioned: face flux from the boundary velocity \\
		Wall, $y = 0$ & $\bb{v}_F = \bb{0}$ & QM: momentum balance; partitioned: flux-consistent gradient \\
		Interface & from the coupling & from the coupling \\
		Solid top & $\bb{D}_S = \bb{0}$ & --- \\
		Solid ends & full: $\bb{D}_S = \bb{0}$; flat: exact $\bb{D}_S$ & --- \\
		\bottomrule
	\end{tabular}
\end{table}

\paragraph{Mesh motion.}
The exact solution does not depend on the motion of the fluid mesh, but the discrete error does.
The interior fluid mesh moves by a displacement Laplacian with quadratic inverse-distance diffusivity.
In the full variant the inlet and outlet points are fixed, because $\hat\eta$ vanishes at the ends and the M\"obius map fixes the end points of the interface; the cells next to the interface then move almost vertically and become non-orthogonal to the tilting interface: $|\bb{k}|/(\hat{\bb{n}}\cdot\bb{d})$ reaches $0.24$ in the interface cells, about $13^\circ$, where $\bb{d}$ is here the vector from the cell centre to the face centre and $\bb{k} = (\bb{I} - \hat{\bb{n}}\hat{\bb{n}})\cdot\bb{d}$.
A stress-based mesh motion halves this to about $6.5^\circ$, with similar errors.
In the flat variant the inlet and outlet points move with the exact linear stretching $y \mapsto y\,h(t)/H$ (with a slip condition there instead, the corner point lags behind the interface and folds the corner cell).
This removes the non-orthogonality at the interface; the inlet and outlet cells, where the prescribed end-point stretching meets the Laplacian motion of the interior, are slightly non-orthogonal ($|\bb{k}|/(\hat{\bb{n}}\cdot\bb{d}) \approx 0.15$).

\paragraph{Error measures.}
The errors are taken against the exact fields at the cell centres of the current mesh (fluid) and of the initial mesh (solid).
They are reported as volume-weighted $L_2$ norms of the pointwise error magnitude, $\|e\|_2 = (\sum_i V_i |e_i|^2/\sum_i V_i)^{1/2}$, and as maxima.
On the interface, the norms of the solid displacement and of the traction applied to the solid are weighted by the face areas of the initial solid interface; the applied traction is compared with $\bb{\sigma}_S\cdot\hat{\bb{n}}_S$, evaluated with the exact normal.
The pressure level is fixed by the outlet, so the pressure error is not corrected by its mean.
The meshes are one cell thick, with $n \times n/2$ fluid and $n \times n/4$ solid square cells for $n$ faces on the interface.
In a temporal study on a fixed mesh, successive time-step solutions are differenced as root-mean-square values over the cells; because the fluid meshes of different time steps differ by the mesh motion, the exact solution is first subtracted at each run's own cell centres.

\paragraph{Checks of the construction.}
The fields are built in sympy and differentiated in symengine, and eleven properties are checked, symbolically or at $40$ significant digits at random points.
The inverse map, $\Xi(x^\Gamma(X,t),t) = X$, and the identity $\bb{\nabla}\cdot(\bb{\nabla}\times\psi) = 0$ for any $\psi$ are checked symbolically.
At the random points the checks are that the top of the fluid domain is the deformed interface, $h(x^\Gamma) = y^\Gamma$; that $\bb{\nabla}\cdot\bb{v}_F = 0$; the kinematic condition and no slip on $y = 0$; traction continuity, with the normal from Nanson's formula; zero solid interface shear for $m_1 = 0$, and the coefficient $c_1$; and the clamps and the zero initial state.
The largest residuals are $3\times10^{-38}$ for the kinematic condition and $1.4\times10^{-38}$ for the traction (both variants), against an interface traction of order one.

\paragraph{Checks of the generated code.}
The closed-form expressions are long, and are evaluated in double precision by generated C++ (common-subexpression elimination; about 3200 temporaries for the fluid fields).
A separate program checks the generated code independently of the symbolic pipeline, at 200 random points.
Fourth-order finite differences of the generated velocity and pressure reproduce $\bb{b}_F$ to $5\times10^{-10}$ and give $\bb{\nabla}\cdot\bb{v}_F$ to $4\times10^{-11}$.
An independent implementation of the neo-Hookean law, written line by line from the source of the code and applied to finite differences of $\bb{D}_S$, reproduces $\bb{\sigma}_S$ to $4\times10^{-10}$ and $\rho_{S,0}\bb{b}_S$ to $2\times10^{-7}$, the truncation error of the nested differences, against components of up to $8$.
The generated solid velocity equals the finite-difference rate of the displacement; the kinematic condition and the interface normal hold; and the fluid and solid tractions agree to $7\times10^{-15}$.
At four reference points the generated values agree with the $40$-digit symbolic values to $2\times10^{-14}$, measured as $|a - b|/(1 + |b|)$.
The program exits with a non-zero status if any check fails.
All checks pass for both variants, and for a set with $\rho_F = 2$ and $\rho_{S,0} = 3$ that tests the density scaling (\texttt{paper-data/mms\_fsi/derive/}).
